\section{Discussion}
\label{sec:discussion}

Cooperative seasonal allocation may reduce volatility relative to independent profit-maximizing
behavior when spatial diversification and sharing rules align member incentives.
Regional and national benchmarks primarily illustrate upper bounds on diversification rather
than practical policy prescriptions.

Welfare comparisons show that sharing rules matter for inequality metrics: equal and
land-weighted rules collapse dispersion in the reference run, while individual retention
and compensation schemes preserve or amplify dispersion depending on deviation from individually
optimal crops.

External validity is limited to the UK arable case study, reference cost and price series,
and cooperative adoption constraints not modeled here.
Forecasting results remain appendix material unless they improve scenario generation.
